\section{Law \& Economics}

\subsection{Accidents: Care \& Activity Levels}
Injurer chooses care $x$ (cost $x$) and activity level $z$; expected harm $z \cdot p(x)H$, $p' < 0$. Efficiency requires optimizing \textbf{both margins}: care where $-p'(x)H = 1$, activity where private benefit $=$ full expected harm.

\textbf{Hand rule} (negligence standard): liable iff $B < PL$ --- untaken precaution cost less than probability $\times$ loss. This is marginal cost-benefit analysis as a legal test.

\subsection{Liability Rules}
\begin{itemize}[nosep]
  \item \textbf{No liability:} victim bears harm $\implies$ victim takes efficient care/activity, injurer takes none
  \item \textbf{Strict liability:} injurer pays all harm $\implies$ injurer efficient on \textit{both} care and activity; victim over-engages (fully insured)
  \item \textbf{Negligence:} injurer liable only if care $<$ standard $x^*$. Injurer meets the standard (care efficient) but then bears no harm at the margin $\implies$ \textbf{activity level too high}. Victim, residual bearer, gets activity right
\end{itemize}
\textbf{Rule of thumb:} the rule assigns the \textit{residual harm}; whoever bears it gets the activity margin right. Choose strict liability when the injurer's activity is the important margin (blasting, wild animals); negligence when the victim's is (both parties' behavior matters).

\textit{TFU trap:} ``negligence and strict liability both induce efficient care, so they are equivalent'' --- F: identical on care, different on activity levels and on who bears residual risk (insurance implication).

\subsection{Property vs.\ Liability Rules (Calabresi--Melamed)}
\textbf{Property rule:} entitlement can only be taken with consent (injunction) --- efficient when transaction costs are low (parties bargain, Coase).

\textbf{Liability rule:} taking allowed at court-set damages --- efficient when transaction costs are high or holdout looms (eminent domain, accidental harms), at the cost of valuation error.

\textit{Link to Coase:} with zero transaction costs the assignment doesn't affect efficiency, only distribution; law matters exactly where bargaining fails---choose the rule that economizes on the relevant friction.

\subsection{Contracts \& Breach}
\textbf{Efficient breach:} breach when performance cost exceeds the promisee's value. \textbf{Expectation damages} (make promisee whole) price breach at its social cost $\implies$ efficient breach decisions; reliance damages underprice breach; specific performance forces renegotiation.

\textit{TFU trap:} penalties above expectation damages deter \textit{efficient} breaches --- overdeterrence is a real cost, not a transfer.
